\documentclass[10pt,a4paper]{article}

\usepackage[utf8]{inputenc}
\usepackage[T1]{fontenc} 

\usepackage{amsmath,mathtools,amsfonts,amssymb}
\usepackage{graphicx}
\usepackage[spanish]{babel}
\usepackage{lastpage,xcolor}
\usepackage[a4paper, top=0.9in, bottom=1in, left=0.9in, right=0.9in]{geometry}
\usepackage{fancyhdr}
\usepackage{enumitem}
\usepackage{titling}
\usepackage[mode=image]{standalone}
\usepackage{tikz}
\usepackage{pgfplots}
\pgfplotsset{compat=1.15}
\usepackage{caption}

\usetikzlibrary{arrows.meta}
\definecolor{utnblue}{RGB}{0, 51, 102}

% Fecha de versión automática según fecha de compilación
\newcommand{\verdate}{\the\year.\ifnum\month<10 0\fi\the\month.\ifnum\day<10 0\fi\the\day}

% Incisos con letras: a), b), c), ...
\setlist[enumerate,1]{label=\textbf{\alph*)}, leftmargin=*, itemsep=4pt}

\setlength{\headsep}{12pt}
\setlength{\droptitle}{-0.5in}
\pretitle{\begin{center}\vspace{-1em}\normalfont\Large}
	\posttitle{\par\end{center}\vspace{-2.5em}}
\preauthor{\begin{center}\vspace{-0.5em}}
	\postauthor{\par\end{center}}
\predate{}\postdate{}

\pagestyle{fancy}
\lhead{\footnotesize Análisis de Señales y Sistemas  - Año \the\year{} Ver. \verdate}
\rhead{\footnotesize Resolución del Trabajo Práctico N$^\text{o}$5}
\rfoot{\footnotesize Página \thepage\ de \pageref{LastPage}}
\renewcommand{\headrulewidth}{0.4pt}
\renewcommand{\footrulewidth}{0.4pt}

% Encabezado de ejercicio
\newcommand{\ejercicio}[1]{\bigskip\noindent{\large\textbf{Ejercicio #1}}\par\nobreak\smallskip}
% Marco teórico breve
\newcommand{\teoria}[1]{\noindent\textit{Marco teórico.}\enspace #1\par\smallskip}

% Diagrama de polos y ceros mínimo: ejes sigma / j-omega
\tikzset{
  pzaxis/.style={->, >=Latex, gray!60},
  polo/.style={utnblue, thick},
  cero/.style={utnblue, thick},
}

\title{%
	{\normalsize Departamento de Ingeniería en Tecnologías Electrónicas, UTN FRM}\\[0.1cm]
	{\large Análisis de Señales y Sistemas}\\[0.15cm]
	{\normalsize Resolución del Trabajo Práctico N$^{\text o}$5: Transformada de Laplace}%
}
\author{}
\date{}

\begin{document}
\maketitle
\thispagestyle{fancy}

\noindent
En los Ejercicios 1 a 12, 15 y 16 se trabaja con señales y sistemas causales
y se usa siempre la transformada de Laplace bilateral. Para una señal
derecha ---en particular, causal--- con transformada racional, la región
de convergencia es siempre el semiplano a la derecha del polo de mayor
parte real. Fijar esa convención de antemano es lo que permite invertir sin
ambigüedad: cada polo aporta su rama causal, $A/(s-p)$ se antitransforma
como $A\,e^{pt}u(t)$, y no hace falta decidir nada más en cada caso. Sin la
ROC, la misma expresión racional admitiría varias señales temporales
distintas; con la convención causal fijada, el problema inverso queda con
solución única.

Los Ejercicios 13 y 14 se apartan de este esquema: allí los sistemas
parten con condiciones iniciales no nulas, y se resuelven con la
transformada de Laplace unilateral, que incorpora esas condiciones
directamente en la propiedad de derivación.

% ==========================================================================
\ejercicio{1: cálculo de la transformada por definición}

\teoria{La idea es no volver a la integral de definición: una señal de
duración finita se arma sumando escalones y rampas que empiezan en distintos
instantes, y cada inicio desplazado se traduce con
$x(t-t_0)\overset{\mathcal{L}}{\longleftrightarrow}e^{-st_0}X(s)$ sobre los
pares básicos $u(t)\leftrightarrow 1/s$ y $t\,u(t)\leftrightarrow 1/s^2$.

Un rasgo común a las cuatro: como la integral se extiende sobre un intervalo
acotado, converge para todo $s$ y la ROC es el plano completo. Por eso un
factor $(1-e^{-s\tau})$ en el numerador siempre cancela el polo aparente en
$s=0$.}

Las cuatro señales están representadas en la Figura~\ref{fig:res_senales}.

\begin{figure}[h]
	\centering
	\includestandalone[width=0.6\textwidth]{../figures/fig_tp5_senales/fig_tp5_senales}
	\caption{Señales de duración finita del Ejercicio 1.}
	\label{fig:res_senales}
\end{figure}

\begin{enumerate}
	\item \emph{Pulso rectangular} de amplitud $A$ y duración $\tau$. Lo pensamos
	como un escalón que enciende la señal en $t=0$, menos otro idéntico
	desplazado a $t=\tau$ que la apaga: $x(t)=A\,[\,u(t)-u(t-\tau)\,]$. El primer
	escalón no está desplazado, así que toma directamente su par básico
	$u(t)\leftrightarrow 1/s$. El segundo, $u(t-\tau)$, es ese mismo escalón con
	el argumento corrido a $t=\tau$: para transformarlo aplicamos la 
	\emph{propiedad de desplazamiento temporal} ---atrasar una señal $t_0$
	unidades en el tiempo multiplica su transformada por $e^{-st_0}$---, que con
	$t_0=\tau$ entrega $u(t-\tau)\leftrightarrow e^{-s\tau}/s$. Restando las dos
	transformadas se obtiene
	\[
		X(s)=A\,\frac{1}{s}-A\,\frac{e^{-s\tau}}{s}=\frac{A}{s}\,\bigl(1-e^{-s\tau}\bigr),
		\qquad \operatorname{ROC}:\ \mathbb{C}.
	\]
	El $1/s$ haría sospechar un polo en el origen, pero el factor $1-e^{-s\tau}$
	se anula exactamente allí ---en $s=0$ vale $1-1=0$---, así que cancela ese
	polo y no queda ninguno: la ROC es todo el plano, como anticipaba el marco
	teórico.

	\item \emph{Rampa que crece de $0$ a $A$ en $[0,\tau]$ y luego se anula.} La
	recta sube una altura $A$ mientras el tiempo avanza un lapso $\tau$, así que su pendiente es $A/\tau$.
	La armamos por superposición. Encendemos en $t=0$ una rampa de pendiente
	$A/\tau$, $\dfrac{A}{\tau}\,t\,u(t)$, que subiría sin techo. En $t=\tau$ esa rampa ya
	vale $A$; le restamos otra rampa de la misma pendiente iniciada en $\tau$,
	$\dfrac{A}{\tau}\,(t-\tau)\,u(t-\tau)$, que cancela la subida y deja la señal plana en
	$A$. Un escalón de altura $A$ iniciado en $\tau$, $-A\,u(t-\tau)$, baja ese tramo
	a cero:
	\[
		x(t)=\frac{A}{\tau}\,t\,u(t)-\frac{A}{\tau}\,(t-\tau)\,u(t-\tau)-A\,u(t-\tau).
	\]
	Cada pieza quedó en función de su propio inicio ---la primera rampa sin
	desplazar; la segunda y el escalón corridos a $\tau$---, así que el
	desplazamiento temporal se aplica directo: la rampa da $1/s^2$, el escalón
	$1/s$, y cada tramo corrido a $\tau$ agrega su $e^{-s\tau}$, con lo que
	\[
		X(s)=\frac{A}{\tau s^2}\,\bigl(1-e^{-s\tau}\bigr)-\frac{A\,e^{-s\tau}}{s},
		\qquad \operatorname{ROC}:\ \mathbb{C}.
	\]

	\item \emph{Señal trapezoidal.} La señal sube de $0$ a $A$ en $[0,\tau_1]$ con
	pendiente $A/\tau_1$, se mantiene en $A$ hasta $\tau_2$ y cae a $0$. La
	armamos por superposición de tres efectos: una rampa que empieza en $0$ y
	subiría indefinidamente, otra de pendiente opuesta que empieza en $\tau_1$
	y cancela la subida ---dejando el tramo plano en $A$--- y un escalón
	negativo en $\tau_2$ que tira la señal a cero:
	\[
		x(t)=\frac{A}{\tau_1}\,t\,u(t)-\frac{A}{\tau_1}\,(t-\tau_1)\,u(t-\tau_1)
		     -A\,u(t-\tau_2).
	\]
	Cada rampa ya está escrita en función de su propio inicio, así que toma
	su par básico $1/s^2$: la primera empieza en $0$ y no lleva exponencial; la
	segunda, corrida a $\tau_1$, recibe por el desplazamiento temporal el factor
	$e^{-s\tau_1}$; y el escalón final, corrido a $\tau_2$, su $e^{-s\tau_2}$
	sobre $1/s$:
	\[
		X(s)=\frac{A}{\tau_1 s^2}\,\bigl(1-e^{-s\tau_1}\bigr)
		     -\frac{A\,e^{-s\tau_2}}{s},
		\qquad \operatorname{ROC}:\ \mathbb{C}.
	\]

	\item \emph{Rampa decreciente.} La señal parte de $A$ en $t=0$ y baja con pendiente
	$-A/\tau$ hasta anularse en $t=\tau$. La leemos como un escalón de altura
	$A$ del que se descuenta una rampa de pendiente $A/\tau$; esa resta sola
	seguiría cayendo por debajo de cero, así que sumamos una tercera rampa que
	empieza en $\tau$ y deja la señal plana en cero a partir de allí:
	\[
		x(t)=A\,u(t)-\frac{A}{\tau}\,t\,u(t)+\frac{A}{\tau}\,(t-\tau)\,u(t-\tau).
	\]
	El escalón da $A/s$ y las dos rampas, $1/s^2$; la segunda empieza en $0$ y
	va sin exponencial, mientras que la tercera, corrida a $\tau$, recibe por el
	desplazamiento temporal su factor $e^{-s\tau}$, de modo que
	\[
		X(s)=\frac{A}{s}-\frac{A}{\tau s^2}\,\bigl(1-e^{-s\tau}\bigr),
		\qquad \operatorname{ROC}:\ \mathbb{C}.
	\]
\end{enumerate}

% ==========================================================================
\ejercicio{2: propiedades}

\teoria{Cada operación sobre la señal en el tiempo se corresponde con una
manipulación sencilla de $X(s)$, y así no hace falta volver a integrar. Las
tres propiedades que usamos son
\[
	\begin{aligned}
		\text{desplazamiento temporal:}&\quad x(t-t_0)\leftrightarrow e^{-st_0}X(s),\\[2pt]
		\text{desplazamiento en }s:&\quad e^{s_0 t}x(t)\leftrightarrow X(s-s_0),\\[2pt]
		\text{derivación en }s:&\quad -t\,x(t)\leftrightarrow \frac{dX}{ds}.
	\end{aligned}
\]
En cada inciso alcanza con reconocer la señal básica escondida y elegir la
propiedad que la lleva a la forma pedida.}

\begin{enumerate}
	\item $x(t)=e^{-a(t-t_0)}u(t-t_0)$. Reconocemos aquí la exponencial básica
	$g(t)=e^{-at}u(t)$ ---con $G(s)=1/(s+a)$ y ROC
	$\operatorname{Re}\{s\}>-a$--- pero con todo su argumento corrido a $t_0$,
	es decir $g(t-t_0)$. Por desplazamiento temporal solo hay que multiplicar
	$G(s)$ por $e^{-st_0}$:
	\[
		X(s)=e^{-st_0}\,G(s)=\frac{e^{-st_0}}{s+a},
		\qquad \operatorname{Re}\{s\}>-a.
	\]
	El desplazamiento no mueve el polo, así que la ROC es la misma que la de
	$g$.

	\item $x(t)=e^{-at}\cos(\omega_0 t)\,u(t)$. La señal de fondo es el coseno
	causal $\cos(\omega_0 t)u(t)$, cuya transformada es $s/(s^2+\omega_0^2)$, y
	lo que la modifica es el factor $e^{-at}$ que la multiplica. Multiplicar
	por una exponencial en el tiempo es justamente el desplazamiento en $s$:
	con $s_0=-a$ hay que reemplazar $s$ por $s+a$ en toda la transformada,
	\[
		X(s)=\frac{(s+a)}{(s+a)^2+\omega_0^2},
		\qquad \operatorname{Re}\{s\}>-a.
	\]
	El reemplazo desplaza los polos de $\pm j\omega_0$ a $-a\pm j\omega_0$: la
	exponencial los corre hacia la izquierda, lo que en el tiempo es la
	envolvente que amortigua al coseno.

	\item $x(t)=t\,e^{-at}u(t)$. Partimos de nuevo de $g(t)=e^{-at}u(t)$,
	$G(s)=1/(s+a)$, pero ahora la señal aparece multiplicada por $t$.
	Multiplicar por $t$ corresponde a derivar respecto de $s$ con un signo,
	$t\,g(t)\leftrightarrow -G'(s)$, así que
	\[
		X(s)=-\frac{d}{ds}\,(s+a)^{-1}=(s+a)^{-2}=\frac{1}{(s+a)^2},
		\qquad \operatorname{Re}\{s\}>-a.
	\]
	El factor $t$ se manifiesta como un polo doble: cada potencia de $t$
	agrega una multiplicidad al polo.

	\item $x(t)=t\sin(at)\,u(t)$. Mismo mecanismo que el inciso anterior,
	ahora sobre el seno causal $\sin(at)u(t)\leftrightarrow a/(s^2+a^2)$.
	Derivando respecto de $s$ y cambiando el signo,
	\[
		X(s)=-\frac{d}{ds}\,\frac{a}{s^2+a^2}
		    =-a\cdot\frac{-2s}{(s^2+a^2)^2}
		    =\frac{2as}{(s^2+a^2)^2},
		\qquad \operatorname{Re}\{s\}>0.
	\]
	Otra vez el factor $t$ eleva la potencia del denominador: el par de polos
	en $\pm ja$ pasa a ser doble.
\end{enumerate}

% ==========================================================================
\ejercicio{3: convolución como producto}

\teoria{El teorema de convolución establece que
$x(t)*h(t)\overset{\mathcal{L}}{\longleftrightarrow}X(s)\,H(s)$: lo que en el
tiempo es una convolución ---una integral--- en el plano $s$ es un simple
producto. La convolución de las dos exponenciales se reduce entonces a
multiplicar sus transformadas y antitransformar el resultado.

Para antitransformar, los residuos de la descomposición en fracciones
parciales se obtienen por cobertura: se multiplica por el factor del polo,
que se cancela, y se evalúa allí. Si el polo $p$ tiene multiplicidad $m$, la
fórmula se generaliza derivando antes de evaluar: para la descomposición
$A_m/(s-p)^m+\cdots+A_1/(s-p)$,
\[
	A_k=\frac{1}{(m-k)!}\,\frac{d^{\,m-k}}{ds^{\,m-k}}\bigl[(s-p)^m\,X(s)\bigr]\bigg|_{s=p},
	\qquad k=m,m-1,\ldots,1.
\]
Con $k=m$ no hay derivada y se recupera la cobertura simple sobre el
coeficiente de mayor potencia; cada potencia $k$ más baja suma una
derivada.}

Con $X(s)=\dfrac{1}{s+a}$ y $H(s)=\dfrac{1}{s+b}$:

\begin{enumerate}
	\item El teorema de convolución convierte la integral temporal en un
	producto, así que sin calcular ninguna integral
	\[
		Y(s)=X(s)\,H(s)=\frac{1}{s+a}\cdot\frac{1}{s+b}=\frac{1}{(s+a)(s+b)}.
	\]

	\item Con $a\neq b$ el producto tiene dos polos simples, en $-a$ y $-b$.
	Los residuos se obtienen por la \emph{fórmula de cobertura}: se
	multiplica $Y(s)$ por el factor del polo correspondiente, que así se
	cancela, y se evalúa el resultado en ese polo:
	\[
		A=\frac{1}{s+b}\bigg|_{s=-a}=\frac{1}{b-a},
		\qquad
		B=\frac{1}{s+a}\bigg|_{s=-b}=\frac{1}{a-b}=-\frac{1}{b-a}.
	\]
	Los dos residuos resultan opuestos, de modo que
	$Y(s)=\dfrac{1}{b-a}\left[\dfrac{1}{s+a}-\dfrac{1}{s+b}\right]$ y, leyendo
	cada término en su rama causal,
	\[
		y(t)=\frac{1}{b-a}\,\bigl(e^{-at}-e^{-bt}\bigr)\,u(t).
	\]
	La convolución de las dos exponenciales resultó ser su diferencia,
	normalizada por la separación $b-a$ entre los polos.

	\item Con $a=b$ los dos polos se confunden en uno solo, ahora doble:
	$Y(s)=1/(s+a)^2$, con polo $p=-a$ y multiplicidad $m=2$. La
	descomposición en fracciones parciales es
	\[
		Y(s)=\frac{A_1}{s+a}+\frac{A_2}{(s+a)^2}.
	\]
	Con $m=2$, la fórmula generalizada de cobertura da primero el
	coeficiente de mayor potencia, $A_2$, con $k=2$ (sin derivar):
	\[
		A_2=(s+a)^2Y(s)\big|_{s=-a}=1.
	\]
	Con $A_2$ en mano, el caso $k=1$ da $A_1$, derivando antes de evaluar:
	\[
		A_1=\frac{d}{ds}\bigl[(s+a)^2Y(s)\bigr]\bigg|_{s=-a}
		   =\frac{d}{ds}[1]_{s=-a}=0.
	\]
	La descomposición resulta entonces $Y(s)=\dfrac{0}{s+a}+\dfrac{1}{(s+a)^2}
	=\dfrac{1}{(s+a)^2}$: un único polo doble no deja nada para separar, y la
	fórmula de multiplicidad lo confirma en vez de
	asumirlo. Con el par $1/(s+a)^2\leftrightarrow t\,e^{-at}u(t)$ del
	inciso 2.c),
	\[
		y(t)=t\,e^{-at}\,u(t).
	\]
	Conviene mirar qué pasó al comparar con el inciso anterior: las dos
	exponenciales distintas colapsan en un único modo, la exponencial
	$e^{-at}$ modulada por el factor $t$ ---el rasgo característico del polo
	doble---. Es el mismo par $1/(a+j\omega)^2$ que aparecía al convolucionar
	estas señales en el dominio de la frecuencia; la transformada de Laplace lo
	recupera con un polo doble en $s=-a$.
\end{enumerate}

\begin{figure}[h]
	\centering
	\includestandalone[width=0.6\textwidth]{../figures/fig_tp5_conv/fig_tp5_conv}
	\caption{Convolución $y(t)$ del Ejercicio 3. (a) Con $a\neq b$ ($a=1$,
	$b=2$): diferencia de dos exponenciales. (b) Con $a=b$ ($a=1$): los dos
	polos se funden en uno doble y aparece el término $t\,e^{-at}$, que sube
	hasta un máximo y luego decae.}
	\label{fig:res_conv}
\end{figure}

% ==========================================================================
\ejercicio{4: transformada inversa con polos reales}

\teoria{Para una $X(s)=P(s)/Q(s)$ propia con polos reales distintos, la
descomposición $X(s)=\sum_k A_k/(s-p_k)$ es única. Los residuos se obtienen
por la fórmula de cobertura $A_k=[(s-p_k)X(s)]_{s=p_k}$: se multiplica por el factor
del polo, que así se cancela, y se evalúa allí.

Cada término aporta su exponencial causal $A_k\,e^{p_k t}u(t)$. Dos casos para
tener a mano: un polo en el origen deja una constante (escalón), y un polo
doble agrega un factor $t$, con sus coeficientes por la fórmula generalizada
de cobertura del Ejercicio 3.}

\begin{enumerate}
	\item $F(s)=\dfrac{2s+9}{(s+3)(s+4)}$. El denominador ya está factorizado,
	con polos simples en $-3$ y $-4$, así que planteamos
	$F(s)=\dfrac{A}{s+3}+\dfrac{B}{s+4}$ y sacamos los residuos por cobertura:
	\[
		A=\frac{2s+9}{s+4}\bigg|_{s=-3}=\frac{3}{1}=3,
		\qquad
		B=\frac{2s+9}{s+3}\bigg|_{s=-4}=\frac{1}{-1}=-1.
	\]
	Cada término aporta su exponencial causal,
	\[
		f(t)=\bigl(3\,e^{-3t}-e^{-4t}\bigr)\,u(t),
	\]
	una suma de dos exponenciales que decaen, con ambos polos en el semiplano
	izquierdo.

	\item $F(s)=\dfrac{4s+8}{s^2+4s-5}$. Aquí el denominador viene sin
	factorizar, y ese es el primer paso: $s^2+4s-5=(s+5)(s-1)$, que deja un
	polo en $-5$ y otro en $+1$. Con los factores a la vista, la cobertura da
	\[
		A=\frac{4s+8}{s-1}\bigg|_{s=-5}=\frac{-12}{-6}=2,
		\qquad
		B=\frac{4s+8}{s+5}\bigg|_{s=1}=\frac{12}{6}=2,
	\]
	y entonces
	\[
		f(t)=\bigl(2\,e^{-5t}+2\,e^{t}\bigr)\,u(t).
	\]
	El polo en $s=1$ ---en el semiplano derecho--- da un término creciente
	$e^{t}$: el sistema es causal pero \emph{inestable}. La posición de un
	único polo basta para delatarlo.

	\item $F(s)=\dfrac{12}{s(s+1)(s+4)}=\dfrac{A}{s}+\dfrac{B}{s+1}+\dfrac{C}{s+4}$.
	Son tres polos simples, uno de ellos en el origen. La cobertura se aplica
	igual a los tres:
	\[
		A=\frac{12}{(s+1)(s+4)}\bigg|_{s=0}=\frac{12}{4}=3,\quad
		B=\frac{12}{s(s+4)}\bigg|_{s=-1}=\frac{12}{-3}=-4,\quad
		C=\frac{12}{s(s+1)}\bigg|_{s=-4}=\frac{12}{12}=1,
	\]
	de donde
	\[
		f(t)=\bigl(3-4\,e^{-t}+e^{-4t}\bigr)\,u(t).
	\]
	El polo en el origen aporta la constante $3\,u(t)$ ---un escalón---, y a
	ella se le suman las dos exponenciales que se apagan.

	\item $F(s)=\dfrac{2s}{(s+1)^2(s+2)}$, con un polo doble en $-1$ y uno
	simple en $-2$. La descomposición es
	\[
		F(s)=\frac{A_1}{s+1}+\frac{A_2}{(s+1)^2}+\frac{B}{s+2}.
	\]
	Sobre el polo doble en $-1$, la fórmula generalizada con $m=2$ da
	primero $A_2$, con $k=2$ (sin derivar):
	\[
		A_2=(s+1)^2 F(s)\big|_{s=-1}=\frac{2s}{s+2}\bigg|_{s=-1}=-2.
	\]
	Con $k=1$ se obtiene $A_1$, derivando antes de evaluar:
	\[
		A_1=\frac{d}{ds}\!\left[\frac{2s}{s+2}\right]_{s=-1}
		   =\frac{2(s+2)-2s}{(s+2)^2}\bigg|_{s=-1}=\frac{4}{1}=4.
	\]
	El coeficiente simple $B$ sale por cobertura directa sobre el polo en
	$-2$:
	\[
		B=(s+2)F(s)\big|_{s=-2}=\frac{2s}{(s+1)^2}\bigg|_{s=-2}=-4.
	\]
	Con $1/(s+1)^2\leftrightarrow t\,e^{-t}u(t)$,
	\[
		f(t)=\bigl(4\,e^{-t}-2t\,e^{-t}-4\,e^{-2t}\bigr)\,u(t).
	\]
	El término $t\,e^{-t}$ es el rasgo característico del polo doble.
\end{enumerate}

% ==========================================================================
\ejercicio{5: transformada inversa con polos complejos}

\teoria{Un par de polos complejos conjugados de una señal real aporta una
sinusoide amortiguada, y hay dos métodos equivalentes para llegar a ella.

El \emph{Método 1} ---residuo--- trata el par como dos polos simples: con
$p=-a+j\omega_0$ se calcula el residuo $A=|A|e^{j\varphi}$, y la
antitransformada del par es el coseno único
$2|A|\,e^{-at}\cos(\omega_0 t+\varphi)\,u(t)$.
 
El \emph{Método 2} ---factor cuadrático--- arma el denominador
$(s+a)^2+\omega_0^2$ a partir del par de polos $-a\pm j\omega_0$; después
acomoda el numerador a la forma $B(s+a)+C\,\omega_0$, que la tabla lee como
$e^{-at}[\,B\cos(\omega_0 t)+C\sin(\omega_0 t)\,]u(t)$.

El Método~1 pasa por los residuos complejos y llega a la respuesta como un único
coseno con desfasaje; el Método~2, el más frecuente en la bibliografía, deja la
respuesta real de entrada en términos de coseno y seno. Resolvemos cada inciso
por ambos, en ese orden, y verificamos que coinciden.}

\begin{enumerate}
	\item $F(s)=\dfrac{3(s+1)}{s^2+2s+5}$. El denominador es un cuadrático
	sin raíces reales: aporta un par de polos complejos conjugados. Lo
	resolvemos por los dos métodos.

	\emph{Método 1.} Ubicamos primero los polos con la resolvente,
	\[
		s=\frac{-2\pm\sqrt{4-20}}{2}=\frac{-2\pm\sqrt{-16}}{2}=-1\pm 2j,
	\]
	es decir $p=-1+2j$ y su conjugado $\bar p=-1-2j$, de donde se leen la
	atenuación $a=1$ y la frecuencia $\omega_0=2$. El residuo del polo $p$
	se obtiene multiplicando por su factor $(s-p)$, que se cancela, y
	evaluando lo que queda en $s=p$:
	\[
		A=\frac{3(s+1)}{s-\bar p}\bigg|_{s=p}
	 =\frac{3(-1+2j+1)}{p-\bar p}
	 =\frac{3(2j)}{4j}=\frac32,
	\]
	El residuo resulta real,
	así que su módulo es $|A|=3/2$ y su fase $\varphi=\arg A=0$. La fórmula del par,
	$2|A|\,e^{-at}\cos(\omega_0 t+\varphi)$, entrega
	\[
		f(t)=2\cdot\tfrac32\,e^{-t}\cos(2t)\,u(t)=3\,e^{-t}\cos(2t)\,u(t),
	\]
	un coseno puro amortiguado por $e^{-t}$: la fase nula indica que no hay
	desfasaje, y por eso no aparece el seno.

	\emph{Método 2.} Las raíces del denominador son $-1\pm2j$, con $a=1$ y
	$\omega_0=2$, así que el par fija el denominador $(s+1)^2+2^2$ y la función
	racional es
	\[
		F(s)=\frac{3(s+1)}{(s+1)^2+2^2}.
	\]
	La tabla lee este cuadrático cuando el numerador tiene la forma
	$B(s+1)+C\cdot2$, con el bloque $(s+1)$ para el coseno y la constante para
	el seno. El numerador $3(s+1)$ ya está así, sin término constante, de modo
	que $B=3$ y $C=0$, y la función queda en la forma de tabla
	\[
		F(s)=3\cdot\frac{s+1}{(s+1)^2+2^2}.
	\]
	Esa única fila se antitransforma directo:
	\[
		f(t)=3\,e^{-t}\cos(2t)\,u(t),
	\]
	el mismo resultado. Que $C=0$ es la contracara de la fase nula del
	Método~1: sin desfasaje no aparece el seno.

	\item $F(s)=\dfrac{5s-11}{s^2+4s+13}$. Otro cuadrático irreducible, que
	resolvemos por los dos métodos.

	\emph{Método 1.} Ubicamos los polos con la resolvente,
	\[
		s=\frac{-4\pm\sqrt{16-52}}{2}=\frac{-4\pm\sqrt{-36}}{2}=-2\pm 3j,
	\]
	o sea $p=-2+3j$ y $\bar p=-2-3j$, con $a=2$ y $\omega_0=3$. El residuo en
	$p$ se obtiene por cobertura,
	\[
		A=\frac{5s-11}{s-\bar p}\bigg|_{s=p}
		 =\frac{5(-2+3j)-11}{6j}
		 =\frac{-21+15j}{6j}
		 =\frac52+\frac72\,j,
	\]
	donde, al separar el cociente, $15j/6j$ da la parte real $\tfrac52$ y
	$-21/6j$ la imaginaria $\tfrac72\,j$ (racionalizando por $j$). Esta vez
	el residuo es complejo, así que su módulo y su fase no son inmediatos:
	\[
		|A|=\frac{\sqrt{5^2+7^2}}{2}=\frac{\sqrt{74}}{2},
		\qquad
		\varphi=\arctan\frac{7}{5}\approx 0{,}951\ \text{rad}.
	\]
	La fórmula del par condensa todo en un coseno con desfasaje,
	\[
		f(t)=\sqrt{74}\,e^{-2t}\cos(3t+0{,}951)\,u(t).
	\]

	\emph{Método 2.} Las raíces del denominador son $-2\pm3j$, con $a=2$ y
	$\omega_0=3$, así que el par fija el denominador $(s+2)^2+3^2$ y la función
	racional es
	\[
		F(s)=\frac{5s-11}{(s+2)^2+3^2}.
	\]
	Para leerla en la tabla, el numerador debe tener la forma $B(s+2)+3C$. A
	diferencia del inciso anterior, el numerador $5s-11$ no está escrito así, y
	hay que acomodarlo: forzamos la aparición del bloque $(s+2)$ y dejamos el
	resto como término constante,
	\[
		5s-11=5(s+2)-10-11=5(s+2)-21=5(s+2)+3(-7),
	\]
	donde el $-21$ se escribió como $3\cdot(-7)$ para que la constante quede
	multiplicada por $\omega_0=3$, como pide la tabla. Con $B=5$ y $C=-7$, la
	función se separa en sus dos términos de tabla,
	\[
		F(s)=5\cdot\frac{s+2}{(s+2)^2+3^2}-7\cdot\frac{3}{(s+2)^2+3^2},
	\]
	y cada término se antitransforma por separado:
	\[
		f(t)=e^{-2t}\bigl[\,5\cos(3t)-7\sin(3t)\,\bigr]\,u(t).
	\]
	Es la misma señal que el Método~1 entregó como un solo coseno desfasado,
	ahora repartida entre coseno y seno. La equivalencia se comprueba sin
	recalcular: la parte real del residuo fija la amplitud del coseno y la
	imaginaria la del seno,
	\[
		2\operatorname{Re}\{A\}=5=B, \qquad -2\operatorname{Im}\{A\}=-7=C.
	\]

	\item $F(s)=\dfrac{s+2}{s(s^2+4)}$. Este caso mezcla los dos tipos de
	polo: uno real y simple en el origen, $s=0$, y un par conjugado que
	proviene de $s^2+4=0$. Antitransformamos la contribución de cada polo por
	separado y sumamos las señales resultantes.

	El polo del origen es real y simple, así que su residuo se obtiene por
	la misma cobertura del Ejercicio~4:
	\[
		A_0=\frac{s+2}{s^2+4}\bigg|_{s=0}=\frac24=\frac12
		\quad\Longrightarrow\quad \tfrac12\,u(t),
	\]
	un escalón de altura $1/2$. Para el par conjugado mostramos los dos métodos.

	\emph{Método 1.} Los polos del par son las raíces de $s^2+4=0$, esto es
	$s=\pm 2j$, con $a=0$ y $\omega_0=2$; que la parte real sea nula ---el par
	cae justo sobre el eje imaginario--- anticipa una oscilación que ni crece
	ni se apaga. Por residuos, el del par en $p=2j$ es
	\[
		A=\frac{s+2}{s\,(s-\bar p)}\bigg|_{s=2j}
		 =\frac{2+2j}{(2j)(4j)}
		 =\frac{2+2j}{-8}
		 =-\frac{1+j}{4},
	\]
	donde el denominador se armó con el factor del propio polo, $s=2j$, y el
	del conjugado, $s-\bar p=4j$. En forma polar, $|A|=\dfrac{\sqrt2}{4}$ y
	$\varphi=\arg A=-\dfrac{3\pi}{4}$. Como $a=0$, la exponencial desaparece y
	el par queda como un coseno puro con desfasaje:
	\[
		2|A|\cos\!\Bigl(2t-\frac{3\pi}{4}\Bigr)
		=\frac{\sqrt2}{2}\cos\!\Bigl(2t-\frac{3\pi}{4}\Bigr)
		=-\tfrac12\cos(2t)+\tfrac12\sin(2t).
	\]
	Sumando ese escalón de altura $1/2$ a la oscilación se obtiene
	\[
		f(t)=\frac12\bigl[\,1-\cos(2t)+\sin(2t)\,\bigr]\,u(t).
	\]

	\emph{Método 2.} El par de polos es $\pm2j$, con $a=0$ y $\omega_0=2$, de
	modo que fija el denominador $(s+0)^2+2^2=s^2+4$. Planteamos la
	descomposición con el par agrupado en ese factor,
	\[
		F(s)=\frac{1/2}{s}+\frac{B\,s+2C}{s^2+4},
	\]
	y para hallar $B$ y $C$ igualamos los numeradores tras llevar todo a
	común denominador,
	\[
		s+2=\tfrac12(s^2+4)+(B s+2C)\,s.
	\]
	Comparamos coeficientes de cada potencia de $s$:
	\[
		\text{(coef. }s^2)\!:\ 0=\tfrac12+B\Rightarrow B=-\tfrac12,
		\qquad
		\text{(coef. }s)\!:\ 1=2C\Rightarrow C=\tfrac12.
	\]
	Como $a=0$, el factor de amortiguamiento $e^{-at}$ vale $1$ y el par
	aporta sin envolvente
	\[
		B\cos(2t)+C\sin(2t)=-\tfrac12\cos(2t)+\tfrac12\sin(2t),
	\]
	idéntico al del Método~1. Sumándole el escalón del origen se obtiene
	\[
		f(t)=\frac12\bigl[\,1-\cos(2t)+\sin(2t)\,\bigr]\,u(t).
	\]
	El polo en el origen deja la constante $1/2$; el par imaginario, una
	oscilación que ni crece ni se amortigua ($a=0$).
\end{enumerate}

\begin{figure}[h]
	\centering
	\includestandalone[width=0.85\textwidth]{../figures/fig_tp5_polos_complejos/fig_tp5_polos_complejos}
	\caption{Antitransformadas con polos complejos del Ejercicio 5. (a) Inciso a:
	coseno amortiguado $3e^{-t}\cos 2t$, con envolvente $\pm 3e^{-t}$ (par de
	polos en $-1\pm 2j$). (b) Inciso b: coseno y seno amortiguados
	$e^{-2t}[5\cos 3t-7\sin 3t]$, con envolvente $\pm\sqrt{74}\,e^{-2t}$ (par de
	polos en $-2\pm 3j$). (c) Inciso c: oscilación sostenida, que ni crece ni se
	apaga (par sobre el eje imaginario, $\pm 2j$).}
	\label{fig:res_polos_complejos}
\end{figure}

% ==========================================================================
\ejercicio{6: de la EDLCC a la función de transferencia}

\teoria{La propiedad de derivación temporal arrastra en general las
condiciones iniciales. Pero acá el sistema parte en reposo, con condiciones
iniciales nulas, así que esos términos se cancelan y cada derivada $d^k/dt^k$
se reduce, sin nada más, a un factor $s^k$. Con esa hipótesis, la ecuación
diferencial colapsa en la función racional $H(s)=P(s)/Q(s)$, que hereda los
mismos coeficientes de la ecuación, y la relación entrada-salida queda
$Y(s)=H(s)X(s)$: es justamente la hipótesis de reposo la que permite definir
$H(s)$ como ese cociente, sin ella habría que arrastrar también la respuesta
natural del estado inicial.

Sus polos ---las raíces de $Q$--- fijan los modos y, con ellos, la estabilidad:
el sistema es BIBO estable si y solo si todos tienen parte real negativa. La
respuesta al impulso se recupera antitransformando $H(s)$ con la ROC causal.}

Para $\ddot y+5\dot y+6y=\dot x+x$:

\begin{enumerate}
	\item Cada derivada pasa a un factor $s$ del orden correspondiente, sin
	términos de condición inicial porque el sistema parte en reposo, y la
	ecuación se vuelve $(s^2+5s+6)Y=(s+1)X$. La función de transferencia es el
	cociente de esos dos polinomios,
	\[
		H(s)=\frac{P(s)}{Q(s)}=\frac{s+1}{s^2+5s+6}=\frac{s+1}{(s+2)(s+3)},
	\]
	donde factorizamos el denominador para dejar los polos a la vista.

	\item El numerador tiene un cero en $z=-1$ y el denominador, polos en
	$p_1=-2$ y $p_2=-3$; la ganancia es $K=b_1/a_2=1$. Los tres están sobre
	el eje real (Figura~\ref{fig:res_pz_ej6}).
	\begin{figure}[h]
		\centering
		\begin{tikzpicture}[scale=0.85, font=\footnotesize]
			\draw[pzaxis] (-3.8,0)--(0.8,0) node[below right]{$\sigma$};
			\draw[pzaxis] (0,-1.1)--(0,1.1) node[left]{$j\omega$};
			\node[polo] at (-2,0) {$\times$};
			\node[polo] at (-3,0) {$\times$};
			\node[cero] at (-1,0) {$\circ$};
			\node[below=3pt] at (-3,0) {$-3$};
			\node[below=3pt] at (-2,0) {$-2$};
			\node[below=3pt] at (-1,0) {$-1$};
		\end{tikzpicture}
		\caption{Constelación del Ejercicio 6: polos ($\times$) en $-2$ y
		         $-3$, cero ($\circ$) en $-1$.}
		\label{fig:res_pz_ej6}
	\end{figure}

	\item La respuesta al impulso es la antitransformada de $H(s)$. Como hay
	dos polos simples, descomponemos $H(s)=\dfrac{A}{s+2}+\dfrac{B}{s+3}$ y
	sacamos los residuos por cobertura,
	\[
		A=\frac{s+1}{s+3}\bigg|_{s=-2}=\frac{-1}{1}=-1,
		\qquad
		B=\frac{s+1}{s+2}\bigg|_{s=-3}=\frac{-2}{-1}=2,
	\]
	y cada uno aporta su exponencial causal,
	\[
		h(t)=\bigl(-e^{-2t}+2\,e^{-3t}\bigr)\,u(t).
	\]

	\item Ambos polos tienen parte real negativa ($-2$ y $-3$), es decir caen
	en el semiplano izquierdo: el sistema es estable, y eso se refleja en que
	su respuesta al impulso decae a cero.
\end{enumerate}

% ==========================================================================
\ejercicio{7: identificación}

\teoria{Si se conocen la entrada $x(t)$ y la salida $y(t)$ de un sistema en
reposo, la relación $Y(s)=H(s)X(s)$ se da vuelta para despejar la transferencia
como el cociente $H(s)=Y(s)/X(s)$.

De $H(s)$ se obtienen después la respuesta al impulso ---antitransformando--- y la
ecuación diferencial ---leyendo $P(s)$ y $Q(s)$ como los polinomios de $x$ y de
$y$ con sus derivadas---.}

\begin{enumerate}
	\item Transformamos las dos señales: $X(s)=\dfrac{1}{s+3}$ y
	$Y(s)=\dfrac{1}{s+1}-\dfrac{1}{s+2}=\dfrac{1}{(s+1)(s+2)}$. El cociente
	deja la transferencia
	\[
		H(s)=\frac{Y(s)}{X(s)}=\frac{s+3}{(s+1)(s+2)},
	\]
	donde dividir por $X$ ---multiplicar por $s+3$--- convirtió el polo de la
	entrada en un cero de $H$. Por cobertura, $A=\dfrac{s+3}{s+2}\big|_{-1}=2$ y
	$B=\dfrac{s+3}{s+1}\big|_{-2}=-1$, de modo que
	$h(t)=(2e^{-t}-e^{-2t})u(t)$. Finalmente, leyendo numerador y denominador
	de $H(s)=\dfrac{s+3}{s^2+3s+2}$ como los coeficientes de la ecuación,
	\[
		\ddot y+3\dot y+2y=\dot x+3x.
	\]

	\item Misma salida, pero ahora $X(s)=\dfrac{2}{s}$:
	\[
		H(s)=\frac{Y(s)}{X(s)}
		    =\frac{1/[(s+1)(s+2)]}{2/s}
		    =\frac{s}{2\,(s+1)(s+2)}.
	\]
	Descomponemos $\dfrac{s}{(s+1)(s+2)}=\dfrac{-1}{s+1}+\dfrac{2}{s+2}$
	(cobertura: $\,s/(s+2)|_{-1}=-1$, $\,s/(s+1)|_{-2}=2$), con lo que
	\[
		h(t)=\tfrac12\bigl(-e^{-t}+2e^{-2t}\bigr)u(t)
		    =\Bigl(-\tfrac12 e^{-t}+e^{-2t}\Bigr)u(t),
	\]
	y de $H(s)=\dfrac{s}{2s^2+6s+4}$,
	\[
		2\ddot y+6\dot y+4y=\dot x
		\qquad\Longleftrightarrow\qquad
		\ddot y+3\dot y+2y=\tfrac12\,\dot x.
	\]
	Los incisos a) y b) comparten la salida pero no la entrada, y conducen a
	sistemas distintos: la salida sola no identifica al sistema; lo identifica
	la relación entrada--salida.

	\item Con $X(s)=\dfrac{1}{s}$ y el par
	$e^{-at}\sin(\omega_0 t)u(t)\leftrightarrow\omega_0/[(s+a)^2+\omega_0^2]$,
	leyendo $a=2$ y $\omega_0=4$ directo de $y(t)=e^{-2t}\sin(4t)u(t)$,
	$Y(s)=\mathcal{L}\{e^{-2t}\sin(4t)u(t)\}=\dfrac{4}{(s+2)^2+16}
	=\dfrac{4}{s^2+4s+20}$:
	\[
		H(s)=\frac{Y(s)}{X(s)}=\frac{4s}{s^2+4s+20}.
	\]
	El denominador de $H(s)$ es el mismo de $Y(s)$, así que sus polos ya
	están identificados: $p=-2\pm j4$, con $a$ y $\omega_0$ leídos de la señal
	de partida. Para antitransformar $h(t)$ acomodamos el numerador:
	$4s=4(s+2)-8$,
	\[
		H(s)=\frac{4(s+2)}{(s+2)^2+16}-2\cdot\frac{4}{(s+2)^2+16},
	\]
	\[
		h(t)=\bigl[\,4\,e^{-2t}\cos(4t)-2\,e^{-2t}\sin(4t)\,\bigr]u(t).
	\]
	De $H(s)=\dfrac{4s}{s^2+4s+20}$,
	\[
		\ddot y+4\dot y+20y=4\dot x.
	\]
	Los polos están en el semiplano izquierdo: el sistema es estable.
\end{enumerate}

% ==========================================================================
\ejercicio{8: leer el plano a partir de $h(t)$}

\teoria{Transformando $h(t)$ se obtiene $H(s)$, y de su forma factorizada se
leen polos y ceros. La parte real del polo gobierna la envolvente
$e^{-\sigma t}$, y la imaginaria, la frecuencia de oscilación.

Esa lectura ya decide la estabilidad: un par sobre el eje imaginario da una
oscilación que no decae ---marginalmente estable, no BIBO---, mientras que un
par en el semiplano izquierdo da una oscilación que decrese en el tiempo ---estable---.}

\begin{enumerate}
	\item $h(t)=2[\cos(5t)+3\sin(5t)]\,u(t)$. Con
	$\cos(5t)u(t)\leftrightarrow s/(s^2+25)$ y
	$\sin(5t)u(t)\leftrightarrow 5/(s^2+25)$, transformamos término a término:
	el $2\cos(5t)$ da $2s/(s^2+25)$, y el $6\sin(5t)$ ---el factor $2$ por el
	$3$ del enunciado--- da $6\cdot 5/(s^2+25)=30/(s^2+25)$. Sumando ambos términos se obtiene
	\[
		H(s)=\frac{2s}{s^2+25}+\frac{30}{s^2+25}=\frac{2(s+15)}{s^2+25}.
	\]
	Polos en $p=\pm j5$ ---sobre el eje imaginario--- y cero en $z=-15$. La
	parte real y la imaginaria del polo ya estaban a la vista en $h(t)$ antes
	de transformar: sin exponencial que la module, $a=0$, y la frecuencia de
	la oscilación fija $\omega_0=5$.
	De $(s^2+25)Y=(2s+30)X$,
	\[
		\ddot y+25y=2\dot x+30x.
	\]
	Los polos están sobre el eje imaginario: el sistema es marginalmente
	estable (su $h(t)$ oscila sin amortiguarse), y por lo tanto \emph{no} es
	BIBO estable.

	\item $h(t)=2e^{-t}[\cos(5t)+3\sin(5t)]\,u(t)$. Es la señal anterior
	multiplicada por $e^{-t}$; por desplazamiento en $s$ ($s\to s+1$),
	\[
		H(s)=\frac{2(s+1)+30}{(s+1)^2+25}=\frac{2s+32}{s^2+2s+26}.
	\]
	Polos en $p=-1\pm j5$ y cero en $z=-16$. Igual que en el inciso anterior,
	la parte real y la imaginaria del polo se leen directo de $h(t)$: la
	envolvente $e^{-t}$ fija $a=1$ y la frecuencia de la oscilación,
	$\omega_0=5$. De
	$(s^2+2s+26)Y=(2s+32)X$,
	\[
		\ddot y+2\dot y+26y=2\dot x+32x.
	\]
	Ahora los polos tienen parte real $-1<0$: el sistema es estable.

	\medskip
	Las dos respuestas oscilan a la misma frecuencia $\omega_0=5$, pero en
	a) la parte real del polo es nula y la oscilación se sostiene, mientras
	que en b) es $-1$ y la oscilación se apaga. Es la parte real del polo
	---su distancia al eje imaginario--- la coordenada que gobierna esa
	diferencia (Figura~\ref{fig:res_pz_ej8}).
	\begin{figure}[h]
		\centering
		\begin{tikzpicture}[scale=0.7, font=\footnotesize]
			% panel a)
			\begin{scope}
			\draw[pzaxis] (-2.2,0)--(1.2,0) node[below right]{$\sigma$};
			\draw[pzaxis] (0,-1.7)--(0,1.7) node[left]{$j\omega$};
			\node[polo] at (0,1.2) {$\times$};
			\node[polo] at (0,-1.2) {$\times$};
			\node[below=2pt] at (-1.4,1.7) {a)};
			\end{scope}
			% panel b)
			\begin{scope}[xshift=5cm]
			\draw[pzaxis] (-2.2,0)--(1.2,0) node[below right]{$\sigma$};
			\draw[pzaxis] (0,-1.7)--(0,1.7) node[left]{$j\omega$};
			\node[polo] at (-0.9,1.2) {$\times$};
			\node[polo] at (-0.9,-1.2) {$\times$};
			\node[below=2pt] at (-1.4,1.7) {b)};
			\end{scope}
		\end{tikzpicture}
		\caption{Polos del Ejercicio 8. a) Sobre el eje imaginario
		         ($\pm j5$): oscilación sostenida. b) En el semiplano
		         izquierdo ($-1\pm j5$): oscilación amortiguada. (Los ceros,
		         en $-15$ y $-16$, quedan fuera del recuadro.)}
		\label{fig:res_pz_ej8}
	\end{figure}
\end{enumerate}

\begin{figure}[h]
	\centering
	\includestandalone[width=0.6\textwidth]{../figures/fig_tp5_h_oscil/fig_tp5_h_oscil}
	\caption{Respuesta al impulso del Ejercicio 8. (a) Inciso a: oscilación
	sostenida (polos en $\pm j5$). (b) Inciso b: la misma frecuencia amortiguada
	por $e^{-t}$, con envolvente $\pm 2\sqrt{10}\,e^{-t}$ (polos en $-1\pm j5$).}
	\label{fig:res_h_oscil}
\end{figure}

% ==========================================================================
\ejercicio{9: leer el plano a partir de la constelación}

\teoria{Una constelación de polos y ceros con su ganancia $K$ determina la
forma factorizada $H(s)=K\,\prod(s-z_i)/\prod(s-p_k)$. De allí se obtienen la
respuesta al impulso, la ecuación diferencial y la estabilidad, que depende de
que todos los polos estén en el semiplano izquierdo.}

De la figura del enunciado: cero en $z=-1$, polos en $p_1=-2$ y $p_2=-6$,
ganancia $K=4$. Entonces
\[
	H(s)=K\,\frac{s-z}{(s-p_1)(s-p_2)}=\frac{4(s+1)}{(s+2)(s+6)}.
\]

\begin{enumerate}
	\item $\displaystyle H(s)=\frac{4(s+1)}{(s+2)(s+6)}$.

	\item Son dos polos simples, así que la antitransformada se obtiene por
	cobertura sobre $H(s)=\dfrac{A}{s+2}+\dfrac{B}{s+6}$:
	\[
		A=\frac{4(s+1)}{s+6}\bigg|_{s=-2}=\frac{4(-1)}{4}=-1,
		\qquad
		B=\frac{4(s+1)}{s+2}\bigg|_{s=-6}=\frac{4(-5)}{-4}=5,
	\]
	y cada residuo acompaña a su exponencial,
	\[
		h(t)=\bigl(-e^{-2t}+5\,e^{-6t}\bigr)\,u(t).
	\]

	\item Reagrupando $H(s)=\dfrac{4s+4}{s^2+8s+12}$, sus coeficientes son los
	de la ecuación,
	\[
		\ddot y+8\dot y+12y=4\dot x+4x.
	\]

	\item Los polos $-2$ y $-6$ caen en el semiplano izquierdo: el sistema es
	estable. El cero en $-1$ no juega aquí ---la estabilidad la deciden solo
	los polos---; su papel es dar forma a la respuesta, no sostenerla o
	apagarla.
\end{enumerate}

% ==========================================================================
\ejercicio{10: interconexión de sistemas}

\teoria{Como la relación $Y=HX$ es algebraica, componer sistemas es operar con
sus transferencias: la cascada las multiplica, el paralelo las suma y la
realimentación negativa ---camino directo $G$, retorno $F$--- da
$G/(1+GF)$. Hay una diferencia decisiva entre las tres: la serie y el paralelo
reúnen los polos de los subsistemas sin moverlos, mientras que la
realimentación los reubica, porque sus polos son las raíces de $1+GF=0$.}

Con $H_1(s)=\dfrac{1}{s+1}$ y $H_2(s)=\dfrac{1}{s+2}$:

\begin{enumerate}
	\item \emph{Serie.} Las transferencias se multiplican,
	\[
		H_{\text{serie}}(s)=H_1(s)\,H_2(s)=\frac{1}{(s+1)(s+2)}.
	\]
	Los polos son los de los dos subsistemas juntos, $s=-1$ y $s=-2$: la
	cascada los apila sin cambiarlos de lugar. Ambos en el semiplano izquierdo,
	así que el sistema es estable.

	\item \emph{Paralelo.} Las transferencias se suman,
	\[
		H_{\text{paralelo}}(s)=H_1(s)+H_2(s)
		    =\frac{1}{s+1}+\frac{1}{s+2}
		    =\frac{(s+2)+(s+1)}{(s+1)(s+2)}=\frac{2s+3}{(s+1)(s+2)}.
	\]
	El denominador es el mismo de la serie, así que los polos siguen en $s=-1$
	y $s=-2$. Lo que aparece es un cero en $s=-3/2$, a mitad de camino entre los
	polos: el paralelo tampoco mueve los polos, solo agrega un cero.

	\item \emph{Realimentación.} Con $G(s)=\dfrac{1}{s-1}$ en el camino directo y
	una ganancia $F(s)=k$ en el de retorno,
	\[
		H(s)=\frac{G(s)}{1+k\,G(s)}
		    =\frac{\dfrac{1}{s-1}}{1+\dfrac{k}{s-1}}
		    =\frac{1}{(s-1)+k}=\frac{1}{s-(1-k)},
	\]
	multiplicando numerador y denominador por $(s-1)$. El polo a lazo cerrado
	quedó en $s=1-k$: la realimentación lo desplazó desde $s=1$ hacia la
	izquierda una distancia $k$. El sistema es estable cuando $1-k<0$, es decir
	$k>1$. El camino directo era inestable ---polo en $+1$---; con $k>1$ la
	realimentación lo lleva al semiplano izquierdo y lo estabiliza.
\end{enumerate}

Las tres configuraciones dejan ver el mismo contraste: la serie y el paralelo
reúnen los polos de los subsistemas sin moverlos ---quedan en $-1$ y $-2$---,
mientras que solo la realimentación reubica el polo. La
Figura~\ref{fig:res_realim} sigue ese corrimiento: el polo de $G$ parte de
$s=1$, en el semiplano derecho, y la realimentación lo corre hacia la izquierda
hasta cruzar al semiplano izquierdo cuando $k>1$.

\begin{figure}[h]
	\centering
	\includestandalone[width=0.5\textwidth]{../figures/fig_tp5_realim/fig_tp5_realim}
	\caption{Polo a lazo cerrado del Ejercicio 10, en $s=1-k$. Parte de $s=1$
	---en el semiplano derecho, inestable, sombreado--- y la realimentación lo
	corre hacia la izquierda; cruza al semiplano izquierdo y el sistema se
	estabiliza cuando $k>1$.}
	\label{fig:res_realim}
\end{figure}

% ==========================================================================
\ejercicio{11: una misma transferencia, distintas entradas}

\teoria{Para un sistema en reposo, la salida a cualquier entrada se obtiene
de $Y(s)=H(s)X(s)$ y una antitransformación. Cada entrada aporta sus propios
polos a los del sistema, y la forma de $y(t)$ se lee de los polos resultantes
de $Y(s)$.}

Con $H(s)=\dfrac{2}{(s+1)(s+2)}$:

\begin{enumerate}
	\item $x(t)=\delta(t)$, $X(s)=1$. Como el impulso transforma a $1$, la
	salida es la propia transferencia, $Y(s)=H(s)$, y antitransformando por
	cobertura ($2/(s+2)|_{-1}=2$, $2/(s+1)|_{-2}=-2$) se obtiene
	\[
		y(t)=h(t)=2\bigl(e^{-t}-e^{-2t}\bigr)u(t).
	\]
	Excitar con $\delta$ es leer directamente la respuesta al impulso.

	\item $x(t)=u(t)$, $X(s)=1/s$: $Y(s)=\dfrac{2}{s(s+1)(s+2)}$. Por
	cobertura $A=2/[(1)(2)]=1$, $B=2/[(-1)(1)]=-2$, $C=2/[(-2)(-1)]=1$:
	\[
		y(t)=\bigl(1-2e^{-t}+e^{-2t}\bigr)u(t).
	\]
	El término permanente es la constante $1=H(0)$; el resto,
	$-2e^{-t}+e^{-2t}$, es el transitorio que se extingue.

	\item $x(t)=e^{-3t}u(t)$, $X(s)=1/(s+3)$. La entrada agrega su propio polo
	en $-3$, distinto de los del sistema, así que
	$Y(s)=\dfrac{2}{(s+1)(s+2)(s+3)}$ tiene tres polos simples. Por cobertura,
	\[
		A=\frac{2}{(s+2)(s+3)}\bigg|_{-1}=\frac{2}{2}=1,\;
		B=\frac{2}{(s+1)(s+3)}\bigg|_{-2}=\frac{2}{-1}=-2,\;
		C=\frac{2}{(s+1)(s+2)}\bigg|_{-3}=\frac{2}{2}=1,
	\]
	de donde
	\[
		y(t)=\bigl(e^{-t}-2e^{-2t}+e^{-3t}\bigr)u(t).
	\]
	Los dos primeros términos son los modos del sistema; el tercero es el que
	trae la entrada.

	\item $x(t)=e^{-t}u(t)$, $X(s)=1/(s+1)$: el polo de la entrada coincide
	con un polo del sistema, y $Y(s)=\dfrac{2}{(s+1)^2(s+2)}$ tiene un polo
	doble en $-1$. Por las fórmulas del Ejercicio 4.d),
	$A_2=2/(s+2)|_{-1}=2$, $B=2/(s+1)^2|_{-2}=2$,
	$A_1=\frac{d}{ds}[2/(s+2)]_{-1}=-2$:
	\[
		y(t)=\bigl(-2e^{-t}+2t\,e^{-t}+2e^{-2t}\bigr)u(t).
	\]
	Para este ejemplo, el factor $t\,e^{-t}$ es el efecto de la coincidencia entre
	el polo de la entrada y el del sistema.
\end{enumerate}

% ==========================================================================
\ejercicio{12: transitorio y permanente ante una sinusoide}

\teoria{La salida se parte en dos: el \emph{transitorio}, formado por los modos
del sistema, que se extingue, y la \emph{respuesta permanente}, sostenida por la entrada.

Para una entrada sinusoidal de frecuencia $\omega_0$, la respuesta permanente es una
sinusoide de la misma frecuencia, con amplitud $|H(j\omega_0)|$ y fase
$\arg H(j\omega_0)$ ---la respuesta en frecuencia evaluada en $\omega_0$---.}

Con $H(s)=\dfrac{1}{s+1}$ y $x(t)=\cos(t)u(t)$, $X(s)=\dfrac{s}{s^2+1}$,
\[
	Y(s)=X(s)H(s)=\frac{s}{(s+1)(s^2+1)}=\frac{s}{(s+1)(s-j)(s+j)},
\]
con un polo real en $-1$ y un par complejo conjugado en $\pm j$. Planteamos
la descomposición en fracciones parciales sobre los tres polos simples,
\[
	Y(s)=\frac{A}{s+1}+\frac{A_p}{s-j}+\frac{\bar A_p}{s+j},
\]
donde el residuo del polo $-j$ es el conjugado del de $j$, porque $Y(s)$
es real.

\begin{enumerate}
	\item El residuo del polo real se obtiene por cobertura,
	\[
		A=\frac{s}{s^2+1}\bigg|_{s=-1}=\frac{-1}{2}=-\frac12.
	\]
	Para el par complejo $s=\pm j$ ---con $a=0$ y $\omega_0=1$--- lo tratamos
	por el Método~1 del Ejercicio~5: el residuo en $p=j$ se obtiene
	multiplicando por su factor $(s-j)$, que se cancela, y evaluando lo que
	queda en $s=j$,
	\[
		A_p=\frac{s}{(s+1)(s+j)}\bigg|_{s=j}=\frac{j}{(j+1)(2j)}=\frac{1}{2(j+1)}=\frac{1-j}{4}.
	\]
	En forma polar, $|A_p|=\dfrac{\sqrt2}{4}$ y $\varphi=\arg A_p=-\dfrac{\pi}{4}$.
	Como $a=0$, la exponencial desaparece y el par queda como un coseno puro
	con desfasaje, $2|A_p|\cos(t+\varphi)=\tfrac{1}{\sqrt2}\cos(t-\pi/4)$.
	Sumando la exponencial del polo real se obtiene
	\[
		y(t)=\Bigl(-\tfrac12 e^{-t}+\tfrac{1}{\sqrt2}\cos\!\bigl(t-\tfrac{\pi}{4}\bigr)\Bigr)u(t).
	\]

	\item El término $-\tfrac12 e^{-t}$ es el \emph{transitorio} ---el modo
	del sistema, asociado al polo $s=-1$, que se apaga---. El coseno
	$\tfrac{1}{\sqrt2}\cos(t-\pi/4)$ es el \emph{permanente}, que subsiste:
	su amplitud y su fase salen directo del residuo del par complejo,
	$2|A_p|$ y $\varphi$.

	\item La respuesta en frecuencia en $\omega_0=1$ es
	$H(j)=\dfrac{1}{1+j}=\dfrac{1-j}{2}$, con
	\[
		|H(j)|=\frac{1}{\sqrt2},
		\qquad
		\arg H(j)=-\frac{\pi}{4}\ \text{rad}\ (-45^\circ).
	\]
	La entrada $\cos t$ tiene amplitud $1$, así que la respuesta permanente resulta con
	amplitud $|H(j)|=1/\sqrt2$ y desfasaje $\arg H(j)=-\pi/4$, exactamente lo
	obtenido en b).
\end{enumerate}

\begin{figure}[h]
	\centering
	\includestandalone[width=0.6\textwidth]{../figures/fig_tp5_trans_perm/fig_tp5_trans_perm}
	\caption{Salida del Ejercicio 12. En rojo, $y(t)$; en azul punteado, el
	término permanente $\tfrac{1}{\sqrt2}\cos(t-\pi/4)$. El transitorio
	$-\tfrac12 e^{-t}$ aparta a $y(t)$ de ese término al comienzo y, al
	extinguirse, las dos curvas se confunden.}
	\label{fig:res_trans_perm}
\end{figure}

% ==========================================================================
\ejercicio{13: EDLCC con condiciones iniciales}

\teoria{La transformada unilateral incorpora el estado inicial: cada derivada
arrastra sus condiciones,
$\mathcal{L}_u\{\dot y\}=sY(s)-y(0^-)$ y
$\mathcal{L}_u\{\ddot y\}=s^2Y(s)-s\,y(0^-)-\dot y(0^-)$.

Transformar la EDLCC la vuelve algebraica; al despejar $Y(s)$ y antitransformar,
la solución aparece ya separada en la parte que se extingue ---el transitorio---
y la que subsiste ---la respuesta permanente---.}

\begin{enumerate}
	\item $\dot y+3y=u(t)$, $y(0)=2$. Aplicamos la propiedad de derivación
	unilateral a $\dot y$, sustituyendo el dato $y(0)=2$:
	\[
		\mathcal{L}_u\{\dot y\}=sY(s)-y(0)=sY(s)-2.
	\]
	El escalón de la derecha transforma a $1/s$, así que la ecuación
	diferencial se vuelve
	\[
		sY-2+3Y=\frac1s.
	\]
	Agrupamos los términos en $Y$ y pasamos el $-2$ al otro miembro, escribiendo
	$2$ como $2s/s$ para sumar con común denominador:
	\[
		(s+3)\,Y=\frac1s+2=\frac{1+2s}{s}=\frac{2s+1}{s}
		\quad\Longrightarrow\quad
		Y(s)=\frac{2s+1}{s\,(s+3)}.
	\]
	Quedan dos polos simples, uno en el origen y otro en $-3$; planteamos
	$Y(s)=\dfrac{A}{s}+\dfrac{B}{s+3}$ y sacamos los residuos por cobertura,
	multiplicando por el factor del polo, que se cancela, y evaluando allí:
	\[ 
		A=\frac{2s+1}{s+3}\bigg|_{s=0}=\frac{1}{3},
		\qquad
		B=\frac{2s+1}{s}\bigg|_{s=-3}=\frac{-5}{-3}=\frac{5}{3}.
	\]
	El polo en el origen antitransforma a una constante y el de $-3$, a su
	exponencial causal:
	\[
		y(t)=\Bigl(\tfrac13+\tfrac53\,e^{-3t}\Bigr)u(t).
	\]
	Conviene chequear que la condición inicial quedó bien incorporada: en
	$t=0$, $y(0)=\tfrac13+\tfrac53=2$, que coincide con el dato de partida. El
	permanente es la constante $1/3$, el valor al que tiende $y(t)$ una vez
	que se apaga el modo propio del sistema; el transitorio es
	$\tfrac53\,e^{-3t}u(t)$, asociado al polo en $-3$, que se extingue.

	\item $\ddot y+4y=0$, $y(0)=1$, $\dot y(0)=2$. Aplicamos la derivación
	unilateral a $\ddot y$, sustituyendo las dos condiciones iniciales:
	\[
		\mathcal{L}_u\{\ddot y\}=s^2Y(s)-s\,y(0)-\dot y(0)=s^2Y(s)-s-2.
	\]
	No hay entrada ---el miembro derecho de la ecuación es cero---, así que la
	ecuación transformada queda directamente
	\[
		s^2Y-s-2+4Y=0.
	\]
	Agrupamos los términos en $Y$ y pasamos $-s-2$ al otro miembro:
	\[
		(s^2+4)\,Y=s+2
		\quad\Longrightarrow\quad
		Y(s)=\frac{s+2}{s^2+4}.
	\]
	El denominador ya tiene la forma de tabla $s^2+\omega_0^2$ con
	$\omega_0=2$; alcanza con separar el numerador en la parte que acompaña a
	$s$ y la constante,
	\[
		Y(s)=\frac{s}{s^2+4}+\frac{2}{s^2+4},
	\]
	donde el primer término es exactamente el par
	$\cos(\omega_0 t)u(t)\leftrightarrow s/(s^2+\omega_0^2)$ y el segundo,
	con la constante igual a $\omega_0=2$, es exactamente el par
	$\sin(\omega_0 t)u(t)\leftrightarrow \omega_0/(s^2+\omega_0^2)$ ---no hace
	falta acomodar nada, porque el numerador ya vino con el $2$ necesario---.
	Antitransformando término a término,
	\[
		y(t)=\bigl(\cos 2t+\sin 2t\bigr)\,u(t).
	\]
	Chequeamos las dos condiciones iniciales: $y(0)=\cos 0+\sin 0=1$, y
	derivando, $\dot y(t)=-2\sin 2t+2\cos 2t$, con $\dot y(0)=2\cos 0=2$;
	ambas coinciden con los datos. No hay entrada: la salida es íntegramente
	respuesta libre, sostenida por la energía inicial. Como los polos $\pm j2$
	están sobre el eje imaginario, la oscilación no se amortigua: no hay un
	transitorio que se apague, y tampoco hay un permanente en el sentido de
	``resto que deja la entrada'', porque aquí no hay entrada.

	\item $\ddot y+2\dot y+5y=e^{-t}u(t)$, $y(0)=0$, $\dot y(0)=1$. Aplicamos
	la derivación unilateral a las dos derivadas, sustituyendo los datos; como
	$y(0)=0$, el término $-s\,y(0)$ de la segunda derivada se anula y solo
	queda $-\dot y(0)=-1$, y en la primera derivada $-y(0)=0$ tampoco aporta:
	\[
		\mathcal{L}_u\{\ddot y\}=s^2Y-s\,y(0)-\dot y(0)=s^2Y-1,
		\qquad
		\mathcal{L}_u\{\dot y\}=sY-y(0)=sY.
	\]
	Con $e^{-t}u(t)\leftrightarrow 1/(s+1)$, la ecuación transformada es
	\[
		s^2Y-1+2sY+5Y=\frac{1}{s+1}.
	\]
	Agrupamos los términos en $Y$ y pasamos el $-1$ al otro miembro sumando
	$1$ de los dos lados:
	\[
		(s^2+2s+5)\,Y=1+\frac{1}{s+1}=\frac{(s+1)+1}{s+1}=\frac{s+2}{s+1},
	\]
	de donde
	\[
		Y(s)=\frac{s+2}{(s+1)(s^2+2s+5)}.
	\]
	El denominador cuadrático no tiene raíces reales: aporta un par de polos
	complejos conjugados, y lo antitransformamos por el Método~1 del
	Ejercicio~5, tratando el par como dos polos simples.

	Ubicamos primero los polos del par con la resolvente sobre
	$s^2+2s+5=0$,
	\[
		s=\frac{-2\pm\sqrt{4-20}}{2}=\frac{-2\pm\sqrt{-16}}{2}=-1\pm 2j,
	\]
	es decir $p=-1+2j$ y $\bar p=-1-2j$, de donde se leen $a=1$ y
	$\omega_0=2$. Planteamos la descomposición sobre los tres polos simples,
	\[
		Y(s)=\frac{A}{s+1}+\frac{A_p}{s-p}+\frac{\bar A_p}{s-\bar p}.
	\]
	El residuo del polo real sale por cobertura, multiplicando por su factor
	$(s+1)$, que se cancela, y evaluando en $s=-1$:
	\[
		A=\frac{s+2}{s^2+2s+5}\bigg|_{s=-1}=\frac14.
	\]
	El residuo del par, en $p$, se obtiene multiplicando por $(s-p)$, que se
	cancela, y evaluando lo que queda en $s=p$:
	\[
		A_p=\frac{s+2}{(s+1)(s-\bar p)}\bigg|_{s=p}
		   =\frac{p+2}{(p+1)(p-\bar p)}
		   =\frac{1+2j}{(2j)(4j)}
		   =\frac{1+2j}{-8}
		   =-\frac18-\frac14\,j,
	\]
	donde el denominador se armó con el factor del propio polo, $p+1=2j$, y
	el del conjugado, $p-\bar p=4j$. El residuo es complejo, así que su
	módulo y su fase salen de sus partes real e imaginaria:
	\[
		|A_p|=\sqrt{\left(\tfrac18\right)^2+\left(\tfrac14\right)^2}=\frac{\sqrt5}{8},
		\qquad
		\varphi=\arg A_p=\arctan(2)-\pi\approx -2{,}034\ \text{rad}\ (-116{,}57^\circ),
	\]
	donde restamos $\pi$ porque el residuo cae en el tercer cuadrante (parte
	real y parte imaginaria negativas). Con $a=1$ y $\omega_0=2$, la fórmula
	del par $2|A_p|\,e^{-at}\cos(\omega_0 t+\varphi)$ entrega su
	contribución,
	\[
		2\cdot\frac{\sqrt5}{8}\,e^{-t}\cos(2t+\varphi)
		=\frac{\sqrt5}{4}\,e^{-t}\cos(2t+\varphi).
	\]
	Sumando la exponencial del polo real se obtiene
	\[
		y(t)=e^{-t}\left[\frac14+\frac{\sqrt5}{4}\cos(2t+\varphi)\right]u(t),
		\qquad \varphi=\arctan(2)-\pi\ \text{rad}\ (-116{,}57^\circ).
	\]
	Esta es la misma señal que el Método~2 (completar cuadrados) entregaría
	como $e^{-t}[\tfrac14-\tfrac14\cos 2t+\tfrac12\sin 2t]$: al expandir el
	coseno desfasado, $\cos\varphi=-1/\sqrt5$ y $\sin\varphi=-2/\sqrt5$
	---las partes real e imaginaria de $A_p$ normalizadas por
	$|A_p|$--- reconstruyen exactamente esos dos coeficientes. Chequeamos la
	condición inicial: en $t=0$,
	\[
		y(0)=\frac14+\frac{\sqrt5}{4}\cos\varphi
		    =\frac14+\frac{\sqrt5}{4}\left(-\frac{1}{\sqrt5}\right)
		    =\frac14-\frac14=0,
	\]
	que coincide con el dato. Todos los términos decaen con $e^{-t}$: como
	la propia entrada se extingue, la respuesta es toda transitoria y el
	permanente es nulo.
\end{enumerate}

% ==========================================================================
\ejercicio{14: un circuito con energía inicial}

\teoria{La transformada unilateral resuelve el circuito con su condición
inicial de un solo paso.

Sobre la solución conviene distinguir dos cortes que no coinciden. Por su
causa, se separa en la respuesta forzada ---debida a la fuente, con el
capacitor inicialmente descargado--- y la respuesta libre ---debida
únicamente a la energía inicial, con la fuente apagada---. Por su destino
temporal, se separa en el transitorio que se apaga y el régimen permanente
al que tiende la tensión. Los dos cortes usan el mismo material pero lo reparten
distinto: la respuesta forzada trae consigo su propio transitorio, así que
no es igual al régimen permanente ni la libre es igual al transitorio completo.}

Con $R=2\ \text{k}\Omega$ y $C=1\ \text{mF}$, la constante de tiempo es
$\tau=RC=2\ \text{s}$, y con $E=10\ \text{V}$ la ecuación del circuito es
$2\,\dot v_C+v_C=10\,u(t)$, con $v_C(0)=4\ \text{V}$. Transformamos con la
propiedad de derivación unilateral, sustituyendo el dato $v_C(0)=4$:
\[
	2\bigl(sV-4\bigr)+V=\frac{10}{s}
	\quad\Longrightarrow\quad
	(2s+1)\,V=\frac{10}{s}+8.
\]

\begin{enumerate}
	\item Despejamos $V(s)$ y separamos en el término que trae la fuente y
	el que trae la condición inicial:
	\[
		V(s)=\frac{10}{s\,(2s+1)}+\frac{8}{2s+1}.
	\]
	Para aplicar cobertura conviene sacar el $2$ como factor común,
	$2s+1=2(s+\tfrac12)$, con lo que el polo distinto del origen queda a la
	vista en $s=-\tfrac12$. El primer término se reescribe
	\[
		\frac{10}{s\cdot 2(s+\frac12)}=\frac{5}{s\,(s+\frac12)}
		=\frac{A}{s}+\frac{B}{s+\frac12},
	\]
	y los residuos salen multiplicando por el factor del polo, que se
	cancela, y evaluando allí:
	\[
		A=\frac{5}{s+\frac12}\bigg|_{s=0}=10,
		\qquad
		B=\frac{5}{s}\bigg|_{s=-\frac12}=-10.
	\]
	El segundo término de $V(s)$ comparte el mismo polo y el mismo factor
	$2$ en el denominador,
	\[
		\frac{8}{2(s+\frac12)}=\frac{4}{s+\frac12},
	\]
	así que sumando los tres términos sobre los dos polos,
	\[
		V(s)=\frac{10}{s}+\frac{-10+4}{s+\frac12}
		    =\frac{10}{s}-\frac{6}{s+\frac12}.
	\]
	Antitransformando,
	\[
		v_C(t)=\Bigl[\,10-6\,e^{-t/2}\,\Bigr]u(t)\ \text{V}.
	\]
	El exponente es $t/2=t/\tau$ con $\tau=RC=2\ \text{s}$, como corresponde
	al polo en $s=-1/\tau$. Chequeamos la condición inicial:
	$v_C(0)=10-6=4\ \text{V}$, que coincide con el dato.

	\item El término constante $10\ \text{V}$ es el \emph{permanente}: la
	tensión del capacitor tiende a la de la fuente $E=10\ \text{V}$. El
	término $-6\,e^{-t/2}\ \text{V}$ es el \emph{transitorio}, que se apaga
	con constante de tiempo $\tau=2\ \text{s}$; su signo negativo dice que el
	capacitor arranca por debajo del valor final y se carga hacia él, lo que
	es consistente con $v_C(0)=4\ \text{V}<E=10\ \text{V}$. Cuando
	$t\to\infty$, $v_C\to 10\ \text{V}$: el valor inicial no interviene en
	ese límite, solo fija el punto de partida y, con él, el tamaño del
	transitorio que hay que recorrer para llegar a $E$.

	Vale la pena mirar también el otro corte, el que separa por causa: del
	inciso a),
	\[
		v_C(t)=\underbrace{10\bigl(1-e^{-t/2}\bigr)}_{\text{forzada}}
		      \,+\,\underbrace{4\,e^{-t/2}}_{\text{libre}}.
	\]
	Este reparto no coincide con el de permanente y transitorio: la
	respuesta forzada trae consigo su propio transitorio, $-10\,e^{-t/2}$,
	que sumado al $4\,e^{-t/2}$ de la libre reconstruye el transitorio
	total $-6\,e^{-t/2}$. La respuesta libre, en cambio, es transitorio
	puro: no aporta nada al valor final $10\ \text{V}$.
\end{enumerate}

\begin{figure}[h]
	\centering
	\includestandalone[width=0.6\textwidth]{../figures/fig_tp5_rc/fig_tp5_rc}
	\caption{Tensión del capacitor del Ejercicio 14 (ejes en unidades
	normalizadas por $E$ y por $\tau$). Nuestro caso, con
	$v_C(0)=4\ \text{V}<E=10\ \text{V}$, corresponde a la curva creciente
	$V_0<E$: el capacitor se carga desde $4\ \text{V}$ hasta $10\ \text{V}$
	con constante de tiempo $\tau=RC=2\ \text{s}$. La otra curva, $V_0>E$,
	ilustra el caso complementario en que el capacitor arranca por encima
	de $E$ y se descarga hacia el mismo valor final.}
	\label{fig:res_rc}
\end{figure}

% ==========================================================================
\ejercicio{15: el plano $s$ contiene a Fourier}

\teoria{La respuesta en frecuencia es la restricción de la transferencia al eje
imaginario, $H(j\omega)=H(s)\big|_{\sigma=0}$, válida cuando ese eje pertenece a
la ROC ---es decir, cuando el sistema es estable---.

Por eso la transformada de Laplace generaliza a la de Fourier: el diagrama
de polos y ceros en el plano $s$ ---qué modos tiene el sistema, si son
oscilantes, si es estable--- ya deja ver su dinámica completa, mientras que
la respuesta en frecuencia es apenas la traza de esa información sobre el
eje imaginario. Restringir a $s=j\omega$ recupera Fourier; pierde de vista
el resto del plano: la posición de los polos y la ROC que certifica la
estabilidad.}

Para $\ddot y+6\dot y+8y=2x$:

\begin{enumerate}
	\item $\displaystyle H(s)=\frac{2}{s^2+6s+8}=\frac{2}{(s+2)(s+4)}$, con
	polos en $-2$ y $-4$. La ROC causal es $\operatorname{Re}\{s\}>-2$.
	Ambos polos están en el semiplano izquierdo: el sistema es estable y el
	eje imaginario pertenece a la ROC.

	\item Dos polos simples, así que la cobertura sobre
	$H(s)=\dfrac{A}{s+2}+\dfrac{B}{s+4}$ da $A=2/(s+4)|_{-2}=1$ y
	$B=2/(s+2)|_{-4}=-1$, de donde
	\[
		h(t)=\bigl(e^{-2t}-e^{-4t}\bigr)u(t),
	\]
	la diferencia de dos exponenciales que decaen.

	\item Como el eje imaginario está en la ROC, se evalúa
	\[
		H(j\omega)=H(s)\big|_{s=j\omega}
		         =\frac{2}{(j\omega)^2+6j\omega+8}
		         =\frac{2}{(8-\omega^2)+j6\omega},
	\]
	que es exactamente la respuesta en frecuencia obtenida con Fourier.
	Restringir $s=j\omega$ recupera Fourier; sobre el resto del plano,
	Laplace aporta la ubicación de los polos y la ROC ---y con ellas la
	estabilidad---, lo que la vuelve aplicable también a sistemas sin
	respuesta en frecuencia (inestables).

	\item $x(t)=e^{-3t}u(t)$, $X(s)=\dfrac{1}{s+3}$. La entrada suma su polo en
	$-3$ a los dos del sistema, de modo que
	$Y(s)=\dfrac{2}{(s+2)(s+3)(s+4)}$ tiene tres polos simples; por cobertura,
	\[
		A=\frac{2}{(s+3)(s+4)}\bigg|_{-2}=\frac{2}{2}=1,\;
		B=\frac{2}{(s+2)(s+4)}\bigg|_{-3}=\frac{2}{-1}=-2,\;
		C=\frac{2}{(s+2)(s+3)}\bigg|_{-4}=\frac{2}{2}=1,
	\]
	de donde
	\[
		y(t)=\bigl(e^{-2t}-2e^{-3t}+e^{-4t}\bigr)u(t).
	\]
	El término en $e^{-3t}$ es el que aporta la entrada; los otros dos son los
	modos propios del sistema.
\end{enumerate}

% ==========================================================================
\ejercicio{16: integrador}

\teoria{Este ejercicio recorre, sobre un único sistema, toda la cadena del
capítulo.

De la EDLCC a la transferencia; de los polos a la respuesta al impulso y la
estabilidad; de la respuesta al escalón a la lectura transitorio--permanente;
del eje imaginario a la respuesta en frecuencia; y de la realimentación al
desplazamiento de los polos.}

Sistema: $\ddot y+4\dot y+3y=\dot x+2x$.

\begin{enumerate}
	\item Transformamos: $H(s)=\dfrac{s+2}{s^2+4s+3}=\dfrac{s+2}{(s+1)(s+3)}$.
	Cero en $z=-2$, polos en $p_1=-1$ y $p_2=-3$, ganancia $K=1$. Los tres
	puntos están sobre el eje real, con el cero entre los dos polos
	(Figura~\ref{fig:res_pz_ej16}).
	\begin{figure}[h]
		\centering
		\begin{tikzpicture}[scale=0.85, font=\footnotesize]
			\draw[pzaxis] (-3.8,0)--(0.8,0) node[below right]{$\sigma$};
			\draw[pzaxis] (0,-1.1)--(0,1.1) node[left]{$j\omega$};
			\node[polo] at (-1,0) {$\times$};
			\node[polo] at (-3,0) {$\times$};
			\node[cero] at (-2,0) {$\circ$};
			\node[below=3pt] at (-3,0) {$-3$};
			\node[below=3pt] at (-2,0) {$-2$};
			\node[below=3pt] at (-1,0) {$-1$};
		\end{tikzpicture}
		\caption{Constelación del Ejercicio 16: polos en $-1$ y $-3$, cero
		         en $-2$.}
		\label{fig:res_pz_ej16}
	\end{figure}

	\item La ROC causal es $\operatorname{Re}\{s\}>-1$. Ambos polos tienen
	parte real negativa: el sistema es causal y estable.

	\item La respuesta al impulso es la antitransformada de $H(s)$. Con dos
	polos simples basta la cobertura sobre $H(s)=\dfrac{A}{s+1}+\dfrac{B}{s+3}$:
	$A=\dfrac{s+2}{s+3}\big|_{-1}=\tfrac12$ y
	$B=\dfrac{s+2}{s+1}\big|_{-3}=\tfrac12$, de donde
	\[
		h(t)=\tfrac12\bigl(e^{-t}+e^{-3t}\bigr)u(t).
	\]

	\item Respuesta al escalón, $X(s)=1/s$:
	$Y(s)=\dfrac{s+2}{s(s+1)(s+3)}$. Por cobertura,
	$A=\dfrac{s+2}{(s+1)(s+3)}\big|_{0}=\tfrac23$,
	$B=\dfrac{s+2}{s(s+3)}\big|_{-1}=-\tfrac12$,
	$C=\dfrac{s+2}{s(s+1)}\big|_{-3}=-\tfrac16$:
	\[
		y(t)=\Bigl(\tfrac23-\tfrac12 e^{-t}-\tfrac16 e^{-3t}\Bigr)u(t).
	\]
	El término permanente es $2/3=H(0)$; los dos términos exponenciales son el
	transitorio.

	\item Como el sistema es estable, el eje imaginario está en la ROC y
	\[
		H(j\omega)=\frac{j\omega+2}{(j\omega)^2+4j\omega+3}
		         =\frac{2+j\omega}{(3-\omega^2)+j4\omega}.
	\]

	\item Realimentación negativa con $G(s)=H(s)$ y $F(s)=k>0$. Los polos a
	lazo cerrado son las raíces de $1+kG(s)=0$, es decir
	\[
		(s+1)(s+3)+k(s+2)=0
		\quad\Longrightarrow\quad
		s^2+(4+k)s+(3+2k)=0.
	\]
	Para una cuadrática $s^2+bs+c=0$, el discriminante $\Delta=b^2-4c$ es
	justamente la cantidad que queda bajo la raíz de la resolvente ---la
	misma que ya usamos en el Ejercicio~5--- y decide el tipo de raíces sin
	necesidad de resolver: si $\Delta>0$ hay dos raíces reales distintas, si
	$\Delta=0$ una raíz real doble, y si $\Delta<0$ un par complejo
	conjugado. Aquí $b=4+k$ y $c=3+2k$, con lo que
	\[
		\Delta=(4+k)^2-4(3+2k)=k^2+4>0
	\]
	para todo $k$, así que las dos raíces son siempre reales: por más que se
	aumente $k$, los polos \emph{no} se vuelven complejos y la respuesta no
	llega a oscilar. Como además ambos coeficientes son positivos para todo
	$k>0$, el sistema permanece estable; al crecer $k$ los polos reales se
	separan, uno corriéndose hacia la izquierda y el otro acercándose al cero
	en $-2$.
\end{enumerate}

\begin{figure}[h]
	\centering
	\includestandalone[width=0.6\textwidth]{../figures/fig_tp5_lugar_raices/fig_tp5_lugar_raices}
	\caption{Lugar de los polos a lazo cerrado del Ejercicio 16 al crecer $k$.
	El discriminante es positivo para todo $k$, así que ambos polos son siempre
	reales: una rama va de $-1$ hacia el cero en $-2$, la otra de $-3$ hacia
	$-\infty$.}
	\label{fig:res_lugar_raices}
\end{figure}

\end{document}
